\documentclass[../../main.tex]{subfiles}
\begin{document}

{\bf[解]}
直円錐の軸を含む平面で切って，下図のように点，図形量を置く．また
\begin{align}
  \cos t = c, \quad \sin t = s
\end{align}
とおく．ただし$0<t<\pi/2$である．

\begin{figure}[htb]
  \centering
  \begin{tikzpicture}[scale=1.8, >=stealth, every node/.style={font=\small}]

    % --- パラメータ設定 ---
    % 球の半径 a, 底角 t (ここでは見た目に合わせて t = 45° または 40°)
    \def\a{1.5}
    \def\tVal{45}

    \pgfmathsetmacro{\sinT}{sin(\tVal)}
    \pgfmathsetmacro{\cosT}{cos(\tVal)}
    \pgfmathsetmacro{\tanT}{tan(\tVal)}

    % 各頂点・中心の幾何座標
    % O: 底面中心（半球の中心）
    % H: 円錐の頂点 (0, a / cos(t))
    % A: 底面の左端 (-a / sin(t), 0)
    % A_R: 底面の右端 (a / sin(t), 0)
    \coordinate (O) at (0, 0);
    \coordinate (H) at (0, {\a / \cosT});
    \coordinate (A) at ({-\a / \sinT}, 0);
    \coordinate (AR) at ({\a / \sinT}, 0);

    % 半球と母線 AH の接点 T (O から AH への垂足)
    % AH の傾きは tan(t) なので、OT の方向角は 180° - (90° - t) = 90° + t
    \coordinate (T) at ({-\a * \sinT}, {\a * \cosT});

    % --- 1. 二等辺三角形 (円錐の軸断面) ---
    \draw[thick] (A) -- (H) -- (AR) -- cycle;

    % --- 2. 高さ OH (点線) ---
    \draw[thick, dotted] (O) -- (H);

    % --- 3. 半円 (半球の断面) ---
    \draw[thick] (\a, 0) arc[start angle=0, end angle=180, radius=\a];

    % --- 4. 半径 a の点線 (O から接点 T へ) ---
    \draw[thick, dotted] (O) -- (T) node[midway, above right] {$a$};

    % --- 5. 角 t の円弧とラベル ---
    % A を中心とし、線分 AH (角度 t) と 底辺 AO (角度 0°) の間の角
    \draw[thick] ($ (A) + (0.2, 0) $) arc[start angle=0, end angle=\tVal, radius=0.2];
    \node at ($ (A) + ({0.5*\tVal}:0.4) $) {$t$};

    % --- 6. 頂点ラベル ---
    \node[above] at (H) {$H$};
    \node[below left] at (A) {$A$};
    \node[below] at (O) {$O$};

  \end{tikzpicture}
  \caption{直円錐の軸断面と内接半球}
  \label{fig:cone_hemisphere_section}
\end{figure}

上図から，
\begin{align}
  &|AO|=\dfrac{a}{c}&|AH|=\dfrac{a}{cs} \\
  &|OH|=\dfrac{a}{c}
\end{align}
であるから，直円錐の表面積$T$は，
\begin{align}
  T&=\text{底面積}+\text{側面積} \\
  &=\pi|AO|^2+\pi|AO||AH| \\
  &=a^2\pi\left(\dfrac{1}{s^2}+\dfrac{1}{cs^2}\right) \\
  &=a^2\pi\dfrac{1}{c(1-c)}\ge 4a^2\pi &\\
\end{align}
等号成立は$c=1/2$の時，つまり$t=\pi/3$の時である．この時，
\[|OH|=2a\]
となる．$\cdots$(答)
\newpage
\end{document}
